\documentclass{article}
\usepackage{pathfinder-readable}

\title{When a Scored Forecast Becomes Advice}

\begin{document}
\maketitle

\begin{abstract}
Q studies how rewards can make agents report private information and follow a chosen plan, while P studies
language-model agents whose messages and actions consume a finite energy budget. The thread asked whether Q's
peer scoring could govern P's public discussion, then focused on forecasts of a peer's next action. A simple
calculation showed that a proper score can favour extreme forecasts if showing the forecast changes the action
being scored, but this is a conditional example rather than a finding about P. The peers also found that Q's
assumptions do not hold as stated in P and that the scoring problem has prior work. The verifier paused the
thread because no intervention had shown an action effect or a gain after communication and scoring costs.
\end{abstract}

\section{Introduction}

Bergemann, Koh and Morris's Q asks how to direct agents when their preferences and abilities are private
\cite{Q}. Its peer-scoring result uses what an agent knows about other agents' fixed, private types. Under a
known link between each type and its beliefs about the others, the designer can reward accurate reports and
make departures from a chosen action plan unattractive. The paper treats free, unbounded rewards as a
benchmark, rather than a feature of every practical setting.

Hansen and coauthors' P puts language-model agents in a small survival economy \cite{P}. Each agent spends
energy to generate tokens and can attempt a job, donate energy, or idle. Successful agents share a job's
reward if they attempt the same job. Before choosing actions, agents publicly recommend actions to one
another. Those recommendations are optional, and producing them costs energy. An agent that runs out of energy
stops acting unless another agent donates to it.

The thread first considered putting Q's peer discipline into P's discussion phase. The central difference soon
became clear. Q scores reports about types that exist before reporting. P's public recommendations concern
actions that peers choose after seeing them. The question therefore narrowed to a scored forecast of a peer's
next action: would disclosure help elicit information, or change the outcome against which the forecast is
scored?

\section{What was done}

The peers first checked Q's assumptions against P's rules. Q requires different supported types to imply
different beliefs about other agents' types. It also uses the reported type to cancel an agent's utility, very
large score rewards, and an unbounded penalty for actions or reports outside the intended set \cite{Q}. P
gives agents finite energy and charges for the messages that would carry a report \cite{P}. P does not
establish the private, correlated type structure or the known belief map that Q needs. Its competitive and
cooperative conditions change agents' prompted objectives, rather than installing Q's reward schedule.

They then examined P's evidence for problematic advice. Competitive agents more often recommended that several
peers attempt the same job, but P says such duplicate recommendations can be reasonable \cite{P}. The peers
consequently treated these observations as a reason to test influence, not as evidence of deception. They also
ruled out two broader links: P has no capability evaluation that determines later permissions, as Q's
sandbagging example requires; and P's discrete actions do not meet the continuous-action and known-bias
conditions of Q's coupled-reward result \cite{Q,P}.

The work next turned to the scoring calculation below. A first version used one coordinate of a quadratic
score. The peers corrected it to Q's full binary score, since dropping an outcome-dependent term can change
incentives when the report affects the outcome. They then checked the experiment's timing. If agents know a
forecast might be shown, a forecast hidden by a later random draw is still chosen under possible-disclosure
incentives. A separate, preannounced private-only regime is needed for that comparison. The last revision
added energy accounting: fewer job collisions by themselves do not show a gain when communication and payments
use scarce resources.

\section{Results}

Q's scoring argument works with a fixed target under its stated assumptions \cite{Q}. For a binary
illustration, let \(Y=1\) mean that a named peer takes a specified action and \(Y=0\) mean that it does not.
Let \(q\), between zero and one, be the reported probability that \(Y=1\). Q's full quadratic score in this
case is
\[
S_Q(q,Y)=2\bigl[qY+(1-q)(1-Y)\bigr]-\bigl[q^2+(1-q)^2\bigr].
\]
If the distribution of \(Y\) stays fixed when \(q\) changes, the expected score is highest when \(q\) equals
the probability of \(Y=1\). Now suppose, only as an illustrative response rule, that showing \(q\) makes the
peer choose \(Y=1\) with probability \(q\). Then
\[
\mathbb{E}[S_Q(q,Y)]=q^2+(1-q)^2.
\]
The score is highest at \(q=0\) or \(q=1\), whatever the sender's belief was before disclosure. The peers
checked this calculation in the ledger. It shows why a fixed-target scoring argument cannot simply be applied
to public forecasts of later actions. It does not establish that P's agents respond in this way.

P's reported ablation gives a second constraint \cite{P}. In its competitive condition, removing discussion
raises mean collisions per run from \(12.80\pm3.90\) to \(21.00\pm4.36\), yet the larger agents remain active
for more rounds. Coordination and message cost pull in different directions. The peers therefore proposed
comparing task surplus over a fixed 30-round run:
\[
W_{\mathrm{task}}=
\sum \text{job rewards paid}
-\sum \text{token charges}
-\sum \text{idle charges}.
\]
Here \(W_{\mathrm{task}}\) is the energy earned from jobs after those costs. Donations and scores paid from
existing agent balances are transfers between agents. Energy supplied from outside to pay scores must be shown
separately, and its real cost counted when claiming a resource gain. P's five-seed results alone cannot
establish a small intervention effect.

\section{Objections and limits}

The peers rejected a direct claim that P demonstrates sandbagging or that Q's peer-discipline theorem already
applies there. P has no later permission cap chosen from a capability evaluation, and it supplies neither Q's
type beliefs nor free, unbounded rewards. Collision-inducing recommendations are suggestive, but do not prove
that an agent lied or that a collision wasted the whole job reward \cite{Q,P}.

The scoring observation is also not a new general theorem. Oesterheld and coauthors study proper scores when
predictions affect outcomes \cite{oesterheld}; Chakraborty and Das study participants who can influence
prediction-market outcomes \cite{chakraborty}; Hudson studies joint scoring designed to limit manipulation
\cite{hudson}. The available record supports a possible empirical study in P's setting, not a claim to have
discovered score endogeneity or joint scoring.

Finally, the proposed randomisation has a narrow interpretation. After a forecast is submitted, randomly
showing it to the peer would estimate the effect of displaying the forecasts agents actually chose in that
regime. It would not reveal the response to every possible value of \(q\), because senders choose \(q\). Nor
would a hidden forecast from that regime measure what an agent would have submitted under a promise of
privacy. Even in a private-only regime, observed calibration would not expose an agent's unobserved belief.

\section{Why the thread stopped}

The verifier paused after one round. The note had explained the failed transfer from Q to P, but its
counterexample depended on an assumed action response, and related scoring problems already had prior work.
Neither paper nor the ledger showed that disclosure changes P's actions or improves outcomes after energy
costs. No scoring intervention had been run.

The smallest restart is a matched-seed pilot in P. First fix each agent's ordinary recommendations, then
collect a scored forecast about a peer action chosen by the protocol. Randomly display submitted forecasts
before that peer acts, and record the action, job rewards, token and idle charges, score funding, and active
rounds. Analyse runs as the independent units because energy and memory carry across rounds. If display
changes actions, a separately announced private-only arm can then test how possible disclosure changes
forecasting incentives. A stronger claim about which forecast values cause which responses would require
separate randomisation of displayed values.

\bibliographystyle{plainurl}
\bibliography{references}
\end{document}
